\chapter{Sviluppi Futuri}

Il lavoro presentato in questa tesi costituisce le fondamenta formali
e l'infrastruttura architetturale del sistema di analisi
\textit{language-agnostic}. La pipeline $\Phi = G \circ R \circ E$,
il Query Engine Algebrico e il sistema di benchmark automatizzato
hanno dimostrato la validità dell'approccio: l'accuratezza
nell'estrazione delle entità architetturali supera il 97\% in tutti i
linguaggi testati, e le relazioni di ereditarietà, import e
invocazione sono correttamente risolte nella maggior parte degli
scenari.

I risultati sperimentali, tuttavia, hanno anche identificato con
precisione le aree in cui il sistema necessita di estensioni. I
prossimi sviluppi sono organizzati in ordine di priorità, partendo
dalle lacune emerse direttamente dai benchmark e procedendo verso
estensioni a più ampio spettro.

\section{Inferenza di Tipo per Variabili Locali}
La limitazione più impattante emersa dai benchmark è l'incapacità del
sistema di inferire il tipo di variabili locali non annotate
esplicitamente. Questa lacuna è responsabile della totalità dei falsi
negativi sugli archi \texttt{accesses\_field} in C (0/4) e C++ (0/4),
e di parte degli errori in Python e Rust.

Il problema si manifesta quando una variabile locale viene dichiarata
tramite un'espressione di inizializzazione il cui tipo è noto
staticamente, ma non annotato:
\begin{verbatim}
  Rectangle rect = {10, 20};     // C: tipo inferibile dal costrutto
  auto car = Car("Fiat", 120);   // C++: tipo inferibile dal costruttore
  admin = Admin("alice", "root") # Python: tipo inferibile dall'istanziazione
\end{verbatim}

L'estensione proposta consiste nell'introduzione di una fase
intermedia di \textbf{propagazione del tipo} ($P_\tau$) tra la Fase 2a
(Query Building) e la Fase 2b (Query Execution). Questa fase
analizzerebbe le assegnazioni all'interno dei blocchi funzionali
(\texttt{Block}) e, per ciascuna variabile locale il cui tipo è
\texttt{Unresolved}, tenterebbe di derivarlo dal lato destro
dell'assegnamento:
\begin{itemize}
  \item Se il lato destro è un'espressione di costruzione
    (\texttt{Instantiates}), il tipo della variabile è il tipo
    costruito.
  \item Se il lato destro è una chiamata a funzione con tipo di
    ritorno noto, il tipo della variabile è il tipo di ritorno.
  \item Se il lato destro è un accesso a campo con tipo noto, il tipo
    della variabile è il tipo del campo.
\end{itemize}

Questa estensione, pur operando localmente a livello di blocco, avrebbe
un impatto significativo sull'accuratezza degli archi nei linguaggi
procedurali (C, C++) e nei linguaggi a tipizzazione dinamica (Python),
dove le annotazioni di tipo esplicite sono rare o assenti.

\section{Risoluzione di Costrutti Mancanti}
I benchmark hanno evidenziato alcune categorie di costrutti
sintattici che il sistema non estrae ancora come siti di dipendenza.
Queste lacune sono puntuali e non richiedono modifiche
all'architettura della pipeline, ma solo estensioni alle euristiche
della Fase 1 e alla semantica del Query Engine.

\subsection{Risoluzione di \texttt{super()}}
In Python e Java, la chiamata \texttt{super()} all'interno di un
metodo restituisce implicitamente un riferimento alla superclasse del
tipo corrente. Il benchmark Python mostra che le chiamate
\texttt{super().\_\_init\_\_()} e \texttt{super().get\_info()} non
vengono estratte (0/2 sulle chiamate a superclasse).

L'implementazione proposta prevede di trattare \texttt{super()} come
una query algebrica speciale:
$$\texttt{super()} \equiv
\text{Extract}(\text{Find}(\texttt{current\_type}), \texttt{super\_types}[0])$$
In questo modo, la risoluzione si riconduce al meccanismo di MRO già
implementato, senza richiedere nuove primitive nel Query Engine.

\subsection{Espressioni di Cast}
I cast di tipo (\texttt{as} in Rust, \texttt{(Type) expr} in Java e
C++) introducono una dipendenza verso il tipo di destinazione che il
sistema attualmente non cattura (0/1 in Rust, 0/2 in Java). Poiché i
nodi CST corrispondenti (\texttt{type\_cast\_expression} in Java,
\texttt{as\_expression} in Rust) sono identificabili dall'euristica,
l'estensione consiste nell'aggiunta di un nuovo tipo di sito di
dipendenza nell'IR (\texttt{casts\_type}) e nella relativa regola di
estrazione in \texttt{body\_extraction}.

\subsection{Campi delle Varianti Enum con Dati Associati}
In Rust, le varianti di enum con dati associati (ad esempio
\texttt{EnumB::FIRST(EnumA)}) generano campi il cui tipo non viene
correttamente risolto. Il problema è localizzato nell'euristica di
parsing delle tuple variant, che attualmente non effettua un
attraversamento ricorsivo dei campi posizionali. La correzione
richiede l'estensione del parser di \texttt{enum\_variant} per
estrarre i tipi dei campi posizionali e mapparli come
\texttt{TypeRef::Unresolved} nell'IR.

\section{Estensione della Rappresentazione Intermedia}
Oltre alle correzioni puntuali, l'IR dovrà essere estesa per coprire
costrutti avanzati che impattano significativamente sull'architettura
di sistemi reali.

\subsection{Generics e Polimorfismo Parametrico}
Linguaggi come Java, C++ (Template) e Rust fanno un uso massiccio di
tipi generici. Attualmente, un riferimento a \texttt{List<User>}
viene estratto come un riferimento al tipo \texttt{List}, perdendo
l'informazione sul parametro di tipo \texttt{User} --- e quindi la
dipendenza architetturale verso di esso.

Il sistema dei tipi $\mathcal{T}$ dovrà essere esteso per catturare i
parametri di tipo:
$$\tau_{generic} = \langle \mathcal{QN}_{base}, \vec{\tau}_{params} \rangle$$
dove $\vec{\tau}_{params}$ è la sequenza dei parametri di tipo (a loro
volta riferimenti di tipo risolvibili). Saranno inoltre necessari i
vincoli sui parametri (es. \texttt{T: Clone} in Rust,
\texttt{T extends Serializable} in Java) per abilitare la risoluzione
corretta delle chiamate a metodo su tipi parametrici.

\subsection{Decoratori e Annotazioni}
In Python e Java, i decoratori (\texttt{@Entity}, \texttt{@Autowired},
\texttt{@dataclass}) alterano semanticamente le classi o i metodi a
cui sono applicati, spesso iniettando dipendenze architetturali
silenziose verso framework esterni (Spring, Django, SQLAlchemy).
Ignorare queste annotazioni significa perdere archi di dipendenza
cruciali per la comprensione dell'architettura. Il benchmark Java
conferma questa lacuna: la dipendenza
\texttt{AnnotationDependency $\to$ SettingsAnnotation} non viene
estratta (0/1 sugli archi di tipo annotazione).

L'estensione proposta prevede l'aggiunta di un nuovo campo
$\vec{\alpha}_{annotations}$ alle entità \texttt{StructuredType} e
\texttt{Function} dell'IR, con il corrispondente tipo di arco
\texttt{AnnotatedBy} nel grafo delle dipendenze.

\subsection{Campi d'Istanza Dichiarati nel Corpo della Classe}
Il benchmark Java evidenzia che i campi dichiarati come variabili
d'istanza nel corpo della classe (ad esempio
  \texttt{DeclarationVariableInstance.pdao} di tipo
\texttt{ProfileDAO}) non generano un arco \texttt{uses\_type}.
L'euristica attuale estrae i campi dalla firma della classe
(parametri del costruttore, dichiarazioni esplicite), ma non
scansiona le dichiarazioni all'interno del corpo della classe stessa.
L'estensione richiede un attraversamento aggiuntivo del
\texttt{class\_body} durante la Fase 1.

\section{Affinamento e Sostituzione delle Euristiche di Estrazione}
L'approccio basato su euristiche di dispatch sui nodi CST di
Tree-sitter ha dimostrato un'efficacia elevata, ma il riconoscimento
dei costrutti tramite euristiche rimane il punto più delicato e fragile
dell'intera pipeline.

\subsection{Evoluzione dell'Estrazione: Integrazione di srcML}
Poiché l'attuale Fase 1 si affida all'analisi testuale e a euristiche sui
nomi dei nodi, è fisiologicamente esposta a variazioni non previste nelle
grammatiche di Tree-sitter. Uno sviluppo futuro particolarmente promettente
riguarda la parziale o totale sostituzione della prima fase di estrazione
(basata su Tree-sitter) con un parser basato su \textbf{srcML}. 
Come discusso nel Capitolo relativo allo Stato dell'Arte, srcML unifica
il formato sintattico mantenendo intatto il codice originale in XML. 
Adottare srcML come base per la Fase 1 eliminerebbe la necessità di
applicare euristiche *ad hoc* per ogni nuovo costrutto, poiché
il parsing e l'etichettatura sintattica sarebbero già standardizzati 
a monte in un formato XML robusto e consolidato, rendendo l'identificazione 
delle entità architetturali un processo puramente dichiarativo (es. tramite query XPath).

\subsection{Gestione delle Macro}
In C, C++ e Rust, le macro espandono codice che non è visibile nel
CST originale. Attualmente, il sistema può perdere dichiarazioni o
riferimenti generati tramite macro (ad esempio, in Rust, il contenuto
  di \texttt{println!()} viene parzialmente perso durante il token
coalescing). Approcci possibili includono:
\begin{itemize}
  \item \textbf{Pre-espansione selettiva}: per C/C++, invocare il
    preprocessore prima dell'analisi per espandere le macro più
    comuni (\texttt{\#define}, \texttt{\#include}).
  \item \textbf{Pattern matching contestuale}: per Rust, riconoscere
    i pattern noti di macro comuni (\texttt{vec![]},
    \texttt{format!()}) e mapparli direttamente come siti di
    dipendenza.
\end{itemize}

\subsection{Graceful Degradation}
Quando il parser incontra porzioni di codice non supportate o
gravemente malformate, l'analisi globale dovrebbe produrre risultati
parziali anziché fallire. Il sistema attuale gestisce già i nodi
\texttt{ERROR} di Tree-sitter ignorandoli, ma non segnala la
copertura effettiva dell'analisi. L'introduzione di un
\textbf{indicatore di confidenza} per ciascun nodo e arco del grafo
(basato sulla presenza di nodi ERROR nelle vicinanze del CST)
permetterebbe all'utente di valutare l'affidabilità dell'output.

\section{Arricchimento del Grafo e Metriche Architetturali}
Una volta che l'astrazione sarà sufficientemente ricca, lo step
finale riguarderà lo sfruttamento algoritmico del Grafo delle
Dipendenze $\mathcal{G} = (V, E)$. Poiché il grafo è
language-agnostic per costruzione, le metriche calcolate su di esso
saranno automaticamente indipendenti dal linguaggio sorgente.

\subsection{Nuove Semantiche di Dipendenza}
L'insieme $\text{Kind}$ degli archi potrà essere espanso per catturare
relazioni architetturali più sofisticate:
\begin{itemize}
  \item \textbf{Ownership e Lifetime} (Rust): distinguere tra
    riferimenti immutabili (\texttt{\&T}), mutabili (\texttt{\&mut T})
    e possesso (\texttt{T}) permetterebbe di tracciare le politiche di
    ownership come archi specializzati.
  \item \textbf{Dipendenze di Testing}: archi tra le classi di test e
    le unità testate, estratti dalla struttura dei file e dalle
    convenzioni di naming.
  \item \textbf{Mixins e Trait con Implementazioni di Default}: la
    gestione dei Mixins (comuni nei linguaggi dinamici) e delle
    implementazioni di default dei Trait richiede una logica di
    risoluzione MRO più sofisticata per tracciare correttamente
    l'iniezione di comportamento.
\end{itemize}

\subsection{Metriche Architetturali}
Il grafo $\mathcal{G}$ fornisce la base per il calcolo automatizzato
di metriche classiche dell'ingegneria del software:
\begin{itemize}
  \item \textbf{Accoppiamento Afferente ed Efferente} ($C_a$, $C_e$):
    il numero di archi entranti e uscenti da ciascun componente,
    calcolabile direttamente dal grafo.
  \item \textbf{Instabilità} ($I = C_e / (C_a + C_e)$): la
    propensione di un componente a cambiare in risposta a cambiamenti
    esterni.
  \item \textbf{Coesione}: misurabile analizzando la connettività
    interna dei sotto-grafi indotti dai moduli.
  \item \textbf{Rilevamento di Cicli di Dipendenza}: l'identificazione
    di componenti fortemente connesse (SCC) nel grafo diretto
    evidenzia le dipendenze circolari, che sono indicatori noti di
    fragilità architetturale.
\end{itemize}

Il valore aggiunto di queste metriche risiede nella loro
\textit{language-agnosticism}: essendo calcolate sul grafo unificato,
producono risultati direttamente confrontabili tra componenti scritti
in linguaggi diversi all'interno dello stesso sistema.

\section{Ottimizzazione delle Prestazioni e Concorrenza}
Il focus principale di questo progetto di tesi è stato lo sviluppo, la
formalizzazione e la validazione del modello e delle strategie di
risoluzione. Di conseguenza, l'implementazione attuale è essenzialmente
\textit{single-threaded} e non è stata ancora sottoposta a
un'ottimizzazione rigorosa delle prestazioni.

Un naturale sviluppo futuro consisterà nel parallelizzare le fasi
dell'analisi. Poiché l'estrazione (Fase 1) opera indipendentemente su
ciascun file sorgente, è intrinsecamente adatta a un'elaborazione
concorrente tramite thread pool (ad esempio, utilizzando la libreria
\texttt{rayon} in Rust). Anche la Fase 2a (Merge) e la Fase 2b (Resolution)
possono essere parzialmente parallelizzate, suddividendo la risoluzione
delle query per namespace o blocco di file, minimizzando i colli di bottiglia
sui blocchi condivisi e riducendo drasticamente i tempi di esecuzione
su codebase di dimensioni enterprise.

\section{Analisi Multi-Linguaggio e Interoperabilità}
Attualmente, il sistema unifica i grafi estratti da progetti multi-linguaggio,
ma risolve le dipendenze in modo isolato all'interno dei confini di ciascun
linguaggio. Una prospettiva futura di grande interesse è la risoluzione di
\textbf{dipendenze inter-linguaggio}.

Nei sistemi complessi è comune che un modulo in Python chiami una libreria in
C/C++ o che un'applicazione Java utilizzi codice nativo tramite \textbf{JNI}
(Java Native Interface). Estendere il motore di risoluzione per mappare
e ricollegare formalmente le chiamate \textit{Foreign Function Interface} (FFI)
permetterebbe di estrarre un Grafo delle Dipendenze realmente globale,
superando completamente la frammentazione architetturale dei microservizi
o dei sistemi polyglot.

\section{Supporto a Paradigmi di Programmazione Alternativi}
L'attuale validazione empirica e la progettazione del modello si sono
concentrate su linguaggi appartenenti principalmente al paradigma
procedurale o orientato agli oggetti (C-like, Java, Python, Rust).

Tuttavia, l'IR è progettata matematicamente per essere agnostica. Un
importante lavoro futuro sarà quello di estendere ed esercitare il modello
su linguaggi appartenenti a paradigmi nettamente differenti. Ad esempio,
valutare la capacità del sistema di catturare le dipendenze in linguaggi
\textbf{Funzionali} (es. Haskell, Lisp, Clojure) basati su funzioni di
ordine superiore e \textit{closure}, o in linguaggi \textbf{Logici}
(es. Prolog). Ciò richiederà un inevitabile affinamento dell'algebra
delle query per supportare semantiche di risoluzione esotiche o puramente
basate sul pattern matching.

\section{Riepilogo delle Priorità}
La Tabella~\ref{tab:priorita-future} sintetizza le estensioni
proposte, il loro impatto atteso sull'accuratezza e la priorità di
implementazione.

\begin{table}[ht]
  \centering
  \caption{Riepilogo delle estensioni future e impatto atteso.}
  \label{tab:priorita-future}
  \begin{tabular}{l l l}
    \hline
    \textbf{Estensione} & \textbf{Impatto Atteso} & \textbf{Priorità} \\
    \hline
    Inferenza tipo variabili locali & C/C++ archi: 40\%$\to$70\%+ & Alta \\
    Risoluzione \texttt{super()}    & Python archi: 63\%$\to$75\%+ & Alta \\
    Campi enum variant (Rust)       & Rust archi: 69\%$\to$80\%+   & Alta \\
    Ottimizzazione Concorrente      & Riduzione drastica tempi esec. & Alta \\
    Espressioni di cast             & +3 archi (Rust, Java)          & Media \\
    Campi d'istanza (Java)          & Java archi: 82\%$\to$88\%+   & Media \\
    Supporto Generics               & Nuove dipendenze parametriche  & Media \\
    Decoratori / Annotazioni        & Dipendenze verso framework     & Media \\
    Sostituzione estrazione (srcML) & Robustezza vs. euristiche testuali & Media \\
    Dipendenze Inter-linguaggio     & Supporto sistemi polyglot (JNI/FFI) & Media \\
    Gestione Macro                  & Copertura C/C++/Rust           & Bassa \\
    Graceful Degradation            & Robustezza su codice reale     & Bassa \\
    Metriche Architetturali         & Sfruttamento del grafo         & Bassa \\
    Paradigmi Alternativi           & Supporto a Lisp/Prolog/Haskell & Bassa \\
    \hline
  \end{tabular}
\end{table}

Nessuna delle estensioni proposte richiede modifiche all'architettura
della pipeline $\Phi$ o alle primitive dell'algebra delle query. Le
correzioni ad alta priorità operano esclusivamente sulle euristiche
di estrazione (Fase 1) e sull'estensione del motore di risoluzione
(Fase 2b), confermando la solidità del modello formale e la
modularità del design complessivo del sistema.
